\chapter{Utility Functions}
\label{ch:utils}

\newthought{The \yamlkey{Utils} block} lets you define your own helper functions and then call
them inside expressions. Today its main feature is \yamlkey{interpolations\_1D}: turning a table
of numbers into a function you can evaluate anywhere a likelihood or mapping is computed.

\section{Why you would want this}
\label{sec:utils-why}

A lot of \textsc{hep} input comes as tables rather than formulas---an experimental limit curve, a
tabulated efficiency, a cross-section fit, a detector response. \yamlkey{interpolations\_1D}
reads such a table and gives you a function \expr{f(x)} that interpolates between the points, so
you can use it directly in a \yamlkey{LogLikelihood} term.

\section{Defining an interpolation}
\label{sec:utils-shape}

\yamlkey{interpolations\_1D} is a list; each item defines one function. You can supply the data
from a \textsc{csv} file or inline.

\paragraph{From a CSV file}

\begin{yamlwin}
Utils:
  interpolations_1D:
    - name: XenonSD2019
      file: "&J/data/xenon_sd_2019.csv"
      logX: false
      logY: true
      kind: cubic
\end{yamlwin}

\begin{description}[style=nextline, leftmargin=1.2em]
  \item[\yamlkey{name}] the function name you will call in expressions.
  \item[\yamlkey{file}] path to a \textsc{csv} with two columns named \code{x} and \code{y}.
  \item[\yamlkey{logX}, \yamlkey{logY}] optional (default \code{false}); set \code{true} when the
        data really lives in log space along that axis.
  \item[\yamlkey{kind}] optional interpolation method: \code{linear}, \code{quadratic},
        \code{cubic}, or \code{nearest}.
\end{description}

\paragraph{Inline}

For a short curve you can write the points directly, with \yamlkey{x\_values} and
\yamlkey{y\_values} instead of \yamlkey{file}:

\begin{yamlwin}
Utils:
  interpolations_1D:
    - name: EffApprox
      x_values: [100, 200, 400, 800]
      y_values: [0.02, 0.10, 0.32, 0.55]
      logX: false
      logY: false
      kind: linear
\end{yamlwin}

\section{Using it in an expression}
\label{sec:utils-use}

Once declared, the name behaves like any other function (Chapter~\ref{ch:symbols}). For example,
comparing a predicted cross-section against a tabulated limit:

\begin{yamlwin}
Sampling:
  LogLikelihood:
    - name: LogL_DD
      expression: "LogGauss(sigSD_p_pb, XenonSD2019(mN1), 0.1 * XenonSD2019(mN1))"
\end{yamlwin}

\noindent It also works in calculator input actions:

\begin{yamlwin}
actions:
  - type: "Dump"
    variables:
      - {name: "limit", expression: "EffApprox(mass)"}
\end{yamlwin}

\begin{jtip}
Use \yamlkey{logX}/\yamlkey{logY} only when the table is genuinely log-spaced. Prefer a
\textsc{csv} file for long experimental tables, and inline values for small control curves. Keep
the function name short and descriptive---you will type it in every expression that uses it.
\end{jtip}

\section{Summary}
\label{sec:utils-summary}

\yamlkey{Utils.interpolations\_1D} turns tabulated data into named functions you can call from any
expression. Supply the data from a two-column \textsc{csv} or inline, choose an interpolation
\yamlkey{kind}, and use the \yamlkey{name} just like a built-in function. This is the bridge
between experimental tables and your likelihood.
